\section{Related Work}
\label{sec:related_work}

\paragraph{ViT token pruning and merging.}
Token reduction methods for Vision Transformers consistently remove 50--66\% of tokens with less than 0.5 percentage points of ImageNet accuracy loss.
DynamicViT~\cite{rao2021dynamicvit} learns binary masks to identify uninformative tokens,
EViT~\cite{liang2022evit} uses attention scores to fuse or discard low-salience patches,
A-ViT~\cite{yin2022avit} introduces adaptive per-token halting depth,
and Token Merging~\cite{bolya2023tome} applies training-free bipartite matching to merge similar tokens.
Their effectiveness rests on a property specific to classification: spatial layout contributes minimally to category identity, so discarding patches incurs little information loss.
Robotic manipulation inverts this assumption.
Grasping and placement tasks depend on the spatial arrangement of objects, and the same token removal rates may destroy task-critical geometric information that classification safely ignores.
None of these methods has been evaluated in closed-loop VLA control.

\paragraph{Token efficiency in VLMs and VLAs.}
Extending token reduction to vision-language models introduces a new importance signal: the language backbone's attention over image tokens.
FastV~\cite{chen2024fastv} observes that image token attention drops sharply after early layers and prunes accordingly.
PruMerge~\cite{shang2024prumerge} clusters and merges spatially adjacent tokens before they enter the language backbone.
These VLM-level methods transfer naturally to VLAs that share the same backbone architecture.
VLA-specific compression has emerged in parallel:
VLA-IAP~\cite{wang2026vlaiap} aligns pruning with action-prediction importance,
VLA-Cache~\cite{hsu2025vlacache} compresses at the KV-cache level,
EfficientVLA~\cite{yang2025efficientvla} combines pruning with merging,
and VLA-Pruner~\cite{liu2025vlapruner} learns a pruning policy jointly with the action head.
Each method validates on a single VLA architecture and demonstrates bounded performance loss on that model.
No study compares pruning robustness across connector designs, so whether a method validated on one connector type transfers to another remains untested.

\paragraph{Vision connectors and token structure.}
The connector between a ViT encoder and a language backbone determines how spatial information is distributed across the tokens that downstream pruning must operate on.
LLaVA~\cite{liu2023llava} projects each patch independently through an MLP, preserving per-token locality.
BLIP-2~\cite{li2023blip2} uses cross-attention with learned queries to aggregate patch information into a fixed set of output tokens.
Pixel shuffle~\cite{shi2016subpixel}, originally designed for super-resolution, folds spatial neighborhoods into channel dimensions and has been adopted by VLAs such as SmolVLA~\cite{shukor2025smolvla} to compress patch groups into single tokens.
These designs create qualitatively different token dependencies: independent connectors (MLP, linear projection) leave patch information separable, while entangling connectors (pixel-shuffle, cross-attention encoders) distribute spatial content across multiple token representations.
The pruning consequences of this distinction have not been characterized.

Our work provides the first cross-architecture analysis of how connector coupling and task complexity jointly determine VLA token compressibility.
Where prior methods optimize pruning for one model, we identify which architectural properties govern pruning tolerance and introduce a single-pass diagnostic (CV of token norms) that correlates with compressibility without requiring closed-loop evaluation.